%%%%%%%%%%%%%%%%%%%%%%%%%%%%%%%%%%%%%%%%%%%%%%%%%%%%%%%%%%%%%%%
%                                                             %
% Cette version repart de 03_Draft\anatomy_current-short.tex  %
%                                                             %
%%%%%%%%%%%%%%%%%%%%%%%%%%%%%%%%%%%%%%%%%%%%%%%%%%%%%%%%%%%%%%%

% NB: Je supprime toutes les (grosses) parties qui étaient en commentaire dans la version précédente pour qu'on y voit plus clair.

We now go further to understand the underlying patterns of these reallocations. First, we describe origin and destination firms with respect to the universe of firms, to assess their specificities. Second, we quantify the role of between/within sector heterogeneity to assess the importance of inter-sectoral transitions in driving our results.

\paragraph{Characterizing origin and destination firms with respect to the universe of firms.}
\label{sec:rolefirm:hiring_firing}

% FORMER VERSION:
% Laid-off workers tend be reallocated towards lower wage premium, higher productivity firms \textit{compared to the control group}. This pattern is interesting as in the cross-section of firms, high value-added per worker is strongly associated with firm wage premium. In this section, we characterize origin and destination firms not only relative to each other---as is already done implicitly by the event-study results presented above---but instead with respect to the entire population of firms. To that end, we simply regress different variables of interest on destination and origin dummies. 

% Results are presented in Table \ref{tab:firing_hiring_rest}. Interestingly, both origin and destination firms are characterized by higher wage premiums than the average firm (Column (1)). Naturally, as suggested by the event-study results, origin firms have higher wage premium than destination ones. Column (2) shows that destination firms exhibit a higher level of productivity than the average firm in the sample. This is also true when measuring productivity as a residual as in Column (3). When comparing origin and destination firms, we see that despite paying higher wage premium, origin firms are less productive. If one considers value-added per worker as a proxy of surplus per worker (or quasi-rent), this suggests that origin firms are characterized by a higher share of surplus going to workers--- which could be rationalized by workers' bargaining power being larger. 

% Columns (4) to (6) repeat the same specification as previously but including sector fixed-effects. Differences with respect to the rest of the firms in the economy are attenuated. The gap between origin and destination firms however remain fairly stable both in terms of wage premiums and productivity. This suggests differences in sectoral composition explains some of the positive gaps between origin and destination firms and the rest of the economy, but do little to explain the gap between these two categories of firms. Column (7) explores further differences in terms of wage premium, controlling for important determinants of such premium, namely size and productivity. Conditional on such variables, destination firms appear to provide lower wage premium than the rest of firms in the sample. The gap between origin and destination firms however remains very stable. The inclusion of sectoral fixed-effects in column (8) hardly affects the estimated coefficients. So, origin and destination firms stand out respectively as high and low-wage firms with respect to the rest of the sample \textit{conditional} on their size and productivity.

The event-study results show that displaced workers are reallocated toward lower-wage-premium and higher-productivity firms. To better understand how these firms compare to the overall population, we now characterize origin and destination firms relative to the universe of firms. To do so, we regress firm-level outcomes on indicators for whether a firm is an origin or a destination firm, including 2-digit sector fixed effects. The coefficients on these indicators measure deviations from the average firm in the same sector.

Results are presented in Table \ref{tab:firing_hiring_rest}.\footnote{Raw firm-level means for the three groups (origin, destination, and other firms) are reported in Appendix Table~\ref{apptab:firing_hiring_means} as a descriptive complement.} %Columns (1)–(3) report unconditional comparisons. Both origin and destination firms exhibit higher wage premia than the average firm (column 1). Consistent with the event-study results, origin firms have higher wage premia than destination firms. At the same time, destination firms are more productive: while origin firms have slightly higher value added per worker than the average firm (column 2), they are less productive than destination firms when comparing the two groups. This pattern is even clearer \textcolor{red}{when using productivity adjusted for capital intensity (column 3)}. Taken together, these results suggest that origin firms combine relatively high wage premia with lower productivity, consistent with a higher share of surplus accruing to workers.
%Columns (4)–(6) extend this comparison by including sector fixed effects, thereby isolating differences within sectors. While the gaps relative to the average firm are attenuated, the differences between origin and destination firms remain largely unchanged. This indicates that sectoral composition explains part of the differences with respect to the overall firm population, but does not account for the gap between origin and destination firms.
Consistent with the event-study results, column (1) shows that origin firms have higher wage premia than destination firms. At the same time, destination firms have a larger value added per worker (column 2) and are more productive (column 3) than both origin firms and other firms. Column (4) further controls for firm size and value added per worker when explaining wage premia. Conditional on these characteristics, destination firms offer lower wage premia than the average firm, while origin firms remain relatively high-premium employers. The origin--destination gap in wage premium, which is larger in column (4) than in column (1), is therefore not driven by differences in observable characteristics such as sector, size, or productivity, but reflects systematic differences in wage-setting. 

Appendix Table~\ref{apptab:firing_hiring_weighted} reports the same comparisons weighted by firm employment, which is equivalent to running the analysis at the worker level relative to an average worker: the qualitative patterns are unchanged.\footnote{Origin and destination firms also differ slightly in age: at $t^D+5$, destination firms are on average $1.5$ years younger than origin firms \textcolor{orange}{[1.9: Appendix Table~\ref{apptab:firm_age_tenure}, $-1.931$ at $t^D+5$]}, and the within-firm worker-mean tenure differs by less than $0.6$ years across the two groups \textcolor{orange}{[now $-0.444$ years and significant (same table): still ``less than 0.6'', but ``about 0.4 years shorter'' would be more informative]} (Appendix Table~\ref{apptab:firm_age_tenure}). The magnitude of these gaps is too small to substitute for the differences in wage-setting we document, but they are consistent with the literature on business dynamism and the average productivity of entering and exiting firms \citep{foster2008reallocation,haltiwanger2018cyclical}.}




\paragraph{Within versus between sector reallocations.} Some studies have highlighted the role of structural change in explaining losses in wage premia \citep{helm2022}. To assess whether this is the case in our context, %the importance of inter-sectoral transitions in driving our results,
we decompose our three main outcomes $Y_{J(i,t)}$ for firm $J$ in sector $s$ %where individual $i$ is employed in $t$ 
(productivity, labor share and firm wage premia) into a (3-digit) -> \textcolor{red}{2-digit} \textcolor{orange}{[OK: the code uses 2-digit sectors]} sector fixed-effect $\bar{Y}_s$ and a residual $R_{J(i,t)} \equiv Y_{J(i,t)}-\bar{Y}_s$. Results are displayed in lighter colors in Figure \ref{fig:ln_apl_ls}. While between-sector effects appear to play an important role in the short-run (sometimes working in the opposite sign than the overall effect) the overall effect is overwhelmingly driven by the within-sector component after $t^D+3$, especially in the case of wage premia (Figure \ref{fig:firm_fe_within_btw}). This implies that, to the extent that inter-sectoral transitions contribute to the pattern of rising productivity and lower wage premia, this occurs through transitions to firms with higher productivity and lower wages with respect to their sector, rather than through transitions across sectors with systematic differences.





\FloatBarrier


